\documentclass[10pt,a4paper]{amsart}
\usepackage[margin=19mm]{geometry}
\usepackage{amsmath,amssymb,mathtools}
\usepackage[scr=rsfs,cal=euler]{mathalfa}
\usepackage{xfrac}
\usepackage[unicode,hidelinks,bookmarks=false]{hyperref}
\newcommand{\ie}{i.e.\ }
\newcommand{\cf}{cf.\ }
\newcommand{\resp}{respectively\ }
\newcommand{\Subsection}{\subsection}
\providecommand{\nobreakdash}{\nobreak\hbox{-}}
\setlength{\emergencystretch}{2em}
\multlinegap0pt
\usepackage{bm}
\usepackage{bbm}
\usepackage{dsfont}
\usepackage{xspace}
\usepackage{cancel}
\usepackage{mathtools}
\usepackage{amsthm}
\makeatletter
\@ifundefined{thm}{\newtheorem{thm}{Thm}}{}
\@ifundefined{lem}{\newtheorem{lem}[thm]{Lemma}}{}
\makeatother
\title[A complete family of Alexandrov-Fenchel inequalities for con]{A complete family of Alexandrov-Fenchel inequalities for convex capillary hypersurfaces in the half-space}
\author{Yingxiang Hu, Yong Wei, Bo Yang, Tailong Zhou}
\date{Source: arXiv:2209.12479v2 (CC BY 4.0)}
\begin{document}
\maketitle

\noindent\emph{Source and licence.} A complete family of Alexandrov-Fenchel inequalities for convex capillary hypersurfaces in the half-space. Version arXiv:2209.12479v2 (CC BY 4.0), distributed under the Creative Commons Attribution 4.0 International licence, as recorded in the arXiv licence metadata for this exact version and shown on the abstract page https://arxiv.org/abs/2209.12479v2. Full rights notice: licenses/arxiv-2209.12479-CC-BY-4.0.txt.\par\smallskip

\noindent\textit{Editorial scope.} Counted unit: the Lem whose source label is \texttt{s3:tech-lemma} in the article (source0.tex lines 718--724, source label \texttt{s3:tech-lemma}), delivered together with the article's own complete original proof of it (source0.tex lines 725--785, one \texttt{proof} environment of 61 lines). No mathematics is written, repaired or re-derived: statement and proof are the article's own text line for line. Required context retained uncounted: s4.Z (lines 714--716). Every definition, display and label of those lines is retained unchanged. Omitted from this package and not redistributed: 1--713, 717--717, 786--1699. Nothing inside a retained excerpt is omitted or altered: every retained body is delivered exactly as the article's lines are written, so no omission receipt of any kind accompanies this package. Citations: the retained excerpts carry no citation command at all, so no bibliography entry of the article is transcribed and no reference is redistributed. Reference adaptation: none was needed and none was made; every reference inside the delivered unit resolves inside it, so no unresolved reference or citation is produced. Printed slips kept literal: no printed word, symbol, hypothesis or number of the counted statement or of its proof was altered. Layout and class: the article's own class is \texttt{amsart} with packages including enumerate, amsrefs, amsfonts, amsmath, amssymb, amscd, amsthm, bm, cancel, mathrsfs, url, graphicx, hyperref, multicol; the wrapper is the lane's pinned \texttt{amsart} editorial wrapper at 10pt on A4, which restates the theorem-like environments the retained text needs and adds the article's own macro definitions that the retained text needs (none), with extra wrapper packages (bm, bbm, dsfont, xspace, cancel, mathtools). The delivered numbering is the unit's own. Assets: no retained span contains a figure environment, a graphics-inclusion command, a PDF-page inclusion, a table or a TikZ picture; no image file is copied into this package, the package declares no asset dependency and no image file is redistributed with it. Licence: version arXiv:2209.12479v2 of this article is distributed under the Creative Commons Attribution 4.0 International licence, as recorded in the arXiv licence metadata for this exact version and shown on the abstract page https://arxiv.org/abs/2209.12479v2. A full rights notice is delivered at licenses/arxiv-2209.12479-CC-BY-4.0.txt. Characters: every retained excerpt is pure ASCII, so the delivered engine, XeLaTeX with its default Latin Modern text font, reports no missing character.
\begin{equation}\label{s4.Z}
Z(p,\xi)=S(p)(\xi,\xi)
\end{equation}

\begin{lem}\label{s3:tech-lemma}
Suppose that $S_{ij}\geq 0$ on $\Sigma$ and
\begin{equation}\label{s4.lem1}
  (\nabla_\mu S_{ij})\xi^i\xi^j\geq 0,\qquad\mathrm{ on}~ \partial\Sigma
\end{equation}
whenever $S_{ij}\xi^j=0$ for some tangent vector $\xi$. If the function $Z$ defined in \eqref{s4.Z} attains its minimum $Z(p,\xi)=0$ at $(p,\xi)\in T\Sigma$, then $D Z|_{(p,\xi)}=0$ and $D^2Z|_{(p,\xi)}$ is non-negative definite.
\end{lem}

\begin{proof}
We choose coordinates $x^1,\cdots,x^n$ for $\Sigma$ near $p$ such that the connection coefficients $\nabla$ vanish at $p$. Then any vector in $T_q{\Sigma}$ for $q$ near $p$ has the form $\sum_{i=1}^n{\dot{x}^i \partial_i}$, so $T\Sigma$ is described locally by the $2n$ coordinates $x^1,\cdots,x^n$ and $\dot{x}^1,\cdots,\dot{x}^n$. We can rotate the coordinates such that $\xi=\partial_1$. In particular, we have $Z(p,\xi)=0$ and $Z(q,\eta)\geq 0$ for all $q\in\Sigma$ and $\eta\in T_q\Sigma$.

If $p\in\operatorname{int}(\Sigma)$, then the conclusion follows easily. So in the following, we assume that $p\in\partial\Sigma$. For any $v\in T_p \partial\Sigma$, we have
   \begin{equation*}
	0=DZ|_{(p,\xi)}\cdot \left(v,0\right),
	\end{equation*}
	We need to show that $DZ|_{(p,\xi)}\cdot \left(\mu,0\right)=0$ as well, where $\mu\in T_p\Sigma$ is the unit outward normal of $\partial\Sigma$ in $\Sigma$. Since $Z(p,\xi)=S_{ij}\xi^i\xi^j=0$,  by the assumption \eqref{s4.lem1}, we have
	\begin{equation*}
	0\leq \left(\nabla_\mu S_{ij}\right)\xi^i\xi^j=DZ|_{(p,\xi)}\cdot \left(\mu,0\right).
	\end{equation*}
If $DZ|_{(p,\xi)}\cdot \left(\mu,0\right)>0$, we consider the geodesic $\gamma:[0,\delta]\rightarrow \Sigma$ with $\gamma(0)=p$ and $\gamma'(0)=-\mu$. Extend $\xi\in T_p\Sigma$ to a vector field $\xi(s)\in T_{\gamma(s)}\Sigma$ along the geodesic $\gamma$ by parallel transport, that is, $\xi'(s)=\nabla_{\gamma'(s)}\xi(s)=0$ and $\xi(0)=\xi$.  Along the geodesic $\gamma(s)$, we have
\begin{align*}
  \frac{d}{ds}\Big |_{s=0}Z(\gamma(s),\xi(s))=&DZ|_{(p,\xi)}\cdot (\gamma'(0),\xi'(0))\\
  =&-DZ|_{(p,\xi)}\cdot \left(\mu,0\right)<0
\end{align*}
 and hence there exists an interior point $\gamma(s_0)\in\text{int}(\Sigma), 0<s_0<\delta$ such that $Z(\gamma(s_0),\xi(s_0))<Z(\gamma(0),\xi(0))=Z(p,\xi)=0$, which contradicts with $Z\geq 0$. Therefore,
 \begin{equation*}
   DZ|_{(p,\xi)}\cdot (v,0)=0,\quad \forall~v\in T_p\Sigma.
 \end{equation*}
 Since $\xi\in T_p\Sigma$ is an interior point, we also have
  \begin{equation*}
   DZ|_{(p,\xi)}\cdot (0,\eta)=0,\quad \forall~\eta\in T_\xi(T_p\Sigma).
 \end{equation*}
 This means that $DZ|_{(p,\xi)}=0$. In local coordinates, we write this as
	\begin{equation*}
	0=\frac{\partial Z}{\partial x^k}=\frac{\partial}{\partial x^k}S_{11}
	\end{equation*}
and
	\begin{equation*}
	0=\frac{\partial Z}{\partial \dot{x}^k}=2S_{k1}
	\end{equation*}
	for all $k=1,\cdots,n$.

Next, we show that $D^2Z$ is non-negative definite, i.e., for any $a=(a_1,\cdots,a_n)$, $b=(b_1,\cdots,b_n)\in\mathbb{R}^n$, at the point $(p,\xi)$ we have
	\begin{equation*}
	D^2Z|_{(p,\xi)}\left((a,b),(a,b)\right)\geq 0.
	\end{equation*}
	Suppose not, then there exists $(a,b)$ such that
	\begin{equation*}
	D^2Z|_{(p,\xi)}\left((a,b),(a,b)\right)<0.
	\end{equation*}
	Note that for any fixed two vectors $\tilde{a},\tilde{b}\in\mathbb{R}^{n}$, we have
	\begin{align*}
	(D^2Z)|_{(p,\xi)}((0,\xi),(\tilde{a},\tilde{b}))=&2\nabla_{\tilde{a}}S_{11}+2\sum_{k=1}^{n}S_{1k}\tilde{b}_k\\
=&DZ|_{(p,\xi)}\cdot \left(2\tilde{a},\tilde{b}\right)=0.
	\end{align*}
	Hence there exits a constant $k\in\mathbb{R}$ such that $b-k\xi\in T_{\xi}\mathbb{S}^{n-1}$ and
	\begin{equation*}
	D^2Z|_{(p,\xi)}\left((a,b-k\xi),(a,b-k\xi))\right)=D^2Z|_{(p,\xi)}\left((a,b),(a,b)\right)<0
	\end{equation*}
	at the point $(p,\xi)$. Without loss of generality, we assume that $\langle a,\mu\rangle<0$, i.e., $a\in T_p\Sigma$ is pointing inward of $\Sigma$. Denote $UT(\Sigma)$ the unit tangent bundle of $\Sigma$, that is, $UT(\Sigma)=\{(p,\xi)\in T\Sigma, ~|\xi|=1\}$. Let $\tilde{\gamma}$ be a curve in $UT(\Sigma)$  such that $\tilde{\gamma}(s)=(p(s),\xi(s))$ with $\tilde{\gamma}(0)=(p,\xi),\ \tilde{\gamma}'(0)=(a,b-k\xi)$. Then we have
	\begin{equation*}
	\frac{d}{ds}\Big|_{s=0}Z(\tilde{\gamma}(s))=DZ|_{(p,\xi)}\cdot(a,b-k\xi)=0,
	\end{equation*}
	and
	\begin{align*} \frac{d^2}{ds^2}\Big|_{s=0}Z(\tilde{\gamma}(s))=&D^2Z|_{(p,\xi)}(\tilde{\gamma}'(0),\tilde{\gamma}'(0))+DZ|_{(p,\xi)}\cdot\tilde{\gamma}''(0)\\
	=&D^2Z|_{(p,\xi)}\left((a,b-k\xi),(a,b-k\xi))\right)<0,
	\end{align*}
	which contradicts with the minimality of $Z(\tilde{\gamma}(s))$ at $s=0$.
\end{proof}

\end{document}
